% RetinexTapetum: hyperparameter selection and result tables
% Required packages:
% \usepackage{amsmath}
% \usepackage{booktabs}
% \usepackage{threeparttable}
% \usepackage{adjustbox}

% -----------------------------------------------------------------------------
% Configuration-selection rule (insert after the HPO description)
% -----------------------------------------------------------------------------
\paragraph{Configuration selection rule.}
Let $\bar{P}_c$, $\bar{S}_c$, and $\bar{L}_c$ denote the mean validation
PSNR, SSIM, and LPIPS, respectively, of candidate configuration $c$ over seeds
$\{42,123,3407\}$. We first retain configurations whose mean PSNR is within
$\delta_P=0.15$~dB of the best observed value,
\begin{equation}
  \mathcal{C}_{P}
  = \left\{c:\bar{P}_c \geq \max_j \bar{P}_j-0.15\right\}.
  \label{eq:hpo-psnr-tolerance}
\end{equation}
Among these configurations, we retain those whose mean SSIM is within
$\delta_S=0.002$ of the highest SSIM in $\mathcal{C}_{P}$,
\begin{equation}
  \mathcal{C}_{S}
  = \left\{c\in\mathcal{C}_{P}:\bar{S}_c
  \geq \max_{j\in\mathcal{C}_{P}}\bar{S}_j-0.002\right\}.
  \label{eq:hpo-ssim-tolerance}
\end{equation}
The final configuration is the member of $\mathcal{C}_{S}$ with the lowest
mean LPIPS,
\begin{equation}
  c^{\star}=\arg\min_{c\in\mathcal{C}_{S}}\bar{L}_c.
  \label{eq:hpo-lpips-selection}
\end{equation}
This lexicographic tolerance rule avoids an arbitrary weighted sum: PSNR is
treated as the primary fidelity criterion, SSIM resolves configurations with
practically similar PSNR, and LPIPS resolves the remaining perceptually similar
candidates. The exact source-profile baseline is evaluated together with the
two leading HPO candidates, so optimization is retained only when it improves
upon the baseline under the same rule. Within each individual training run,
the stored checkpoint is the epoch with the highest validation PSNR. The
official test partitions are not consulted during configuration, seed, or
checkpoint selection.

% -----------------------------------------------------------------------------
% Table 1: fixed parameters shared by all datasets
% -----------------------------------------------------------------------------
\begin{table}[t]
  \centering
  \caption{Training, architecture, and loss parameters fixed across all
  RetinexTapetum experiments.}
  \label{tab:retinextapetum-fixed-hyperparameters}
  \small
  \begin{threeparttable}
    \begin{adjustbox}{max width=\columnwidth}
      \begin{tabular}{lll}
        \toprule
        Category & Parameter & Fixed value \\
        \midrule
        Training & Maximum epochs & 120 \\
        Training & Batch size & 2 \\
        Training & Crop size & $256\!\times\!256$ \\
        Architecture & Base channels & 128 \\
        Architecture & Trainable parameters & 536,463 \\
        Optimizer & Optimizer & Adam \\
        Optimizer & $(\beta_1,\beta_2)$ & $(0.9,0.999)$ \\
        Optimizer & Learning-rate schedule & Cosine annealing \\
        Enhancement loss & $w_{\mathrm{L1}}$ & 1.00 \\
        Decomposition loss & $w_{\mathrm{recon\text{-}low}}$ & 1.00 \\
        Decomposition loss & $w_{\mathrm{recon\text{-}high}}$ & 0.85 \\
        Decomposition loss & $w_{\mathrm{reflect}}$ & 0.06 \\
        Decomposition loss & $w_{\mathrm{smooth\text{-}low}}$ & 0.08 \\
        Decomposition loss & $w_{\mathrm{smooth\text{-}high}}$ & 0.08 \\
        Enhancement loss & $w_{\mathrm{smooth\text{-}enh}}$ & 0.06 \\
        Final confirmation & Random seeds & $\{42,123,3407\}$ \\
        \bottomrule
      \end{tabular}
    \end{adjustbox}
    \begin{tablenotes}[flushleft]
      \footnotesize
      \item The 120-epoch value is the maximum training horizon. A run may end
      earlier, but only validation-selected checkpoints are used for reporting.
    \end{tablenotes}
  \end{threeparttable}
\end{table}

% -----------------------------------------------------------------------------
% Table 2: selected dataset-specific parameters
% -----------------------------------------------------------------------------
\begin{table*}[t]
  \centering
  \caption{Dataset-specific hyperparameters selected for RetinexTapetum.}
  \label{tab:retinextapetum-hyperparameters}
  \small
  \begin{threeparttable}
    \begin{adjustbox}{max width=\textwidth}
      \begin{tabular}{lcccc}
        \toprule
        Hyperparameter
          & LOL-v1
          & LOL-v2 Real
          & LOL-v2 Synthetic
          & UHD-LL (down4) \\
        \midrule
        Initial learning rate
          & $6.9063\!\times\!10^{-4}$
          & $6.6209\!\times\!10^{-4}$
          & $6.9916\!\times\!10^{-4}$
          & $4.6948\!\times\!10^{-4}$ \\
        Minimum learning rate
          & $1.4628\!\times\!10^{-5}$
          & $4.0830\!\times\!10^{-5}$
          & $5.0238\!\times\!10^{-5}$
          & $2.1283\!\times\!10^{-5}$ \\
        Gradient-clipping norm
          & 3 & 3 & 8 & 3 \\
        $\lambda_{\max}$
          & 1.80 & 2.00 & 1.30 & 1.35 \\
        $w_{\mathrm{SSIM}}$
          & 0.50 & 0.45 & 0.65 & 0.35 \\
        $w_{\mathrm{LPIPS}}$
          & 0.10 & 0.12 & 0.08 & 0.08 \\
        $w_{\mathrm{color}}$
          & 0.04 & 0.02 & 0.14 & 0.06 \\
        $w_{\mathrm{chroma}}$
          & 0.04 & 0.14 & 0.12 & 0.12 \\
        $w_{\mathrm{edge}}$
          & 0.07 & 0.02 & 0.00 & 0.04 \\
        $w_{\mathrm{dark\text{-}noise}}$
          & 0.03 & 0.01 & 0.05 & 0.01 \\
        $w_{\mathrm{attn}}$
          & 0.009 & 0.010 & 0.007 & 0.012 \\
        \bottomrule
      \end{tabular}
    \end{adjustbox}
    \begin{tablenotes}[flushleft]
      \footnotesize
      \item Parameters not listed here are fixed as reported in
      Table~\ref{tab:retinextapetum-fixed-hyperparameters}. The UHD-LL down4
      column corresponds to the source-profile baseline because it remained the
      best configuration after final three-seed confirmation.
    \end{tablenotes}
  \end{threeparttable}
\end{table*}

% -----------------------------------------------------------------------------
% Table 3: final three-seed configuration-selection results
% -----------------------------------------------------------------------------
\begin{table*}[t]
  \centering
  \caption{Three-seed validation performance of the selected configurations.}
  \label{tab:retinextapetum-hpo-three-seed}
  \small
  \begin{threeparttable}
    \begin{adjustbox}{max width=\textwidth}
      \begin{tabular}{llccc}
        \toprule
        Dataset & Selected source
          & PSNR $\uparrow$ & SSIM $\uparrow$ & LPIPS $\downarrow$ \\
        \midrule
        LOL-v1
          & HPO-C2
          & $21.907 \pm 0.077$
          & $0.8652 \pm 0.0022$
          & $0.1256 \pm 0.0055$ \\
        LOL-v2 Real
          & HPO-C2
          & $21.576 \pm 0.037$
          & $0.8317 \pm 0.0004$
          & $0.1615 \pm 0.0059$ \\
        LOL-v2 Synthetic
          & HPO-C1
          & $25.173 \pm 0.045$
          & $0.9357 \pm 0.0014$
          & $0.0554 \pm 0.0011$ \\
        UHD-LL (down4)
          & Source baseline
          & $27.035 \pm 0.178$
          & $0.9263 \pm 0.0021$
          & $0.0670 \pm 0.0027$ \\
        \bottomrule
      \end{tabular}
    \end{adjustbox}
    \begin{tablenotes}[flushleft]
      \footnotesize
      \item Values are mean $\pm$ standard deviation over seeds
      $\{42,123,3407\}$. Configuration selection uses
      Eqs.~\eqref{eq:hpo-psnr-tolerance}--\eqref{eq:hpo-lpips-selection}.
      Final-confirmation LPIPS is evaluated after resizing to 512 pixels.
    \end{tablenotes}
  \end{threeparttable}
\end{table*}

% -----------------------------------------------------------------------------
% Table 4: selected-checkpoint training/validation and independent test results
% -----------------------------------------------------------------------------
\begin{table*}[t]
  \centering
  \caption{Training diagnostics, validation performance, and independent
  test-set performance of the selected RetinexTapetum checkpoints.}
  \label{tab:retinextapetum-train-test}
  \scriptsize
  \setlength{\tabcolsep}{3.2pt}
  \begin{threeparttable}
    \begin{adjustbox}{max width=\textwidth}
      \begin{tabular}{lcccccccccccc}
        \toprule
        & & &
        \multicolumn{2}{c}{Loss at best epoch} &
        \multicolumn{4}{c}{Validation split} &
        \multicolumn{4}{c}{Independent test split} \\
        \cmidrule(lr){4-5}
        \cmidrule(lr){6-9}
        \cmidrule(lr){10-13}
        Dataset
          & Seed
          & Best/completed
          & Train
          & Val.
          & $N$
          & PSNR $\uparrow$
          & SSIM $\uparrow$
          & LPIPS $\downarrow$
          & $N$
          & PSNR $\uparrow$
          & SSIM $\uparrow$
          & LPIPS $\downarrow$ \\
        \midrule
        LOL-v1
          & 42 & 88/120
          & 0.2779 & 0.2713
          & 48 & 22.009 & 0.8681 & 0.1185
          & 15 & 21.039 & 0.8316 & 0.1357 \\
        LOL-v2 Real
          & 123 & 74/119
          & 0.2819 & 0.2864
          & 69 & 21.598 & 0.8318 & 0.1532
          & 100 & 20.019 & 0.8405 & 0.1258 \\
        LOL-v2 Synthetic
          & 3407 & 111/120
          & 0.1780 & 0.1809
          & 90 & 25.117 & 0.9375 & 0.0538
          & 100 & 25.053 & 0.9322 & 0.0577 \\
        UHD-LL (down4)
          & 3407 & 96/120
          & 0.1587 & 0.1554
          & 116 & 27.259 & 0.9281 & 0.0652
          & 150 & 24.711 & 0.9115 & 0.0884 \\
        \bottomrule
      \end{tabular}
    \end{adjustbox}
    \begin{tablenotes}[flushleft]
      \footnotesize
      \item Best/completed reports the selected checkpoint epoch and the number
      of completed training epochs. The training pipeline logs train loss but
      not train-split PSNR, SSIM, or LPIPS. Validation metrics correspond to
      the same seed and checkpoint used for test evaluation. Validation LPIPS
      is evaluated with 512-pixel resizing, whereas test LPIPS is taken from
      the paper-metric evaluation output; validation--test LPIPS differences
      should therefore be interpreted cautiously. The test partitions are used
      only after all HPO and checkpoint-selection decisions are complete.
    \end{tablenotes}
  \end{threeparttable}
\end{table*}
